% =====================================================================
%  COURSE CONTENT  —  page 2 of the certificate
%  This file is included by certificate.tex via \courseContentFile.
%  To use a different course, copy this file, edit the content, and point
%  \courseContentFile (in certificate.tex) at the new file name.
%
%  Available helpers / styles (defined in certificate.tex):
%    \coursesection{Title}   heading with an accent underline
%    colours: ink, softink, accent   (e.g. \color{ink})
%  Standard LaTeX (itemize, tabular, \textbf, ...) works normally and
%  will flow onto further pages if the content is long.
% =====================================================================

\coursesection{Course Overview}
This programme provides a structured awareness of the FSSC~22000 food safety
management scheme, built on ISO~22000:2018 and the ISO/TS~22002 prerequisite
programmes. Participants learn how the scheme is structured, what it requires
of an organisation, and how it is audited and maintained.

\vspace{6pt}
\begin{tabular}{@{}p{3.2cm}p{9cm}@{}}
  \textbf{Duration}   & 8 hours (1 day) \\[2pt]
  \textbf{Delivery}   & Instructor-led — classroom or live online \\[2pt]
  \textbf{Level}      & Awareness / Foundation \\[2pt]
  \textbf{Assessment} & Multiple-choice examination, 70\% pass mark \\[2pt]
  \textbf{Prerequisites} & None \\
\end{tabular}

\coursesection{Learning Objectives}
On completion of this course, participants will be able to:
\begin{itemize}\setlength\itemsep{2pt}
  \item Explain the purpose, structure and benefits of FSSC~22000 (v5.1).
  \item Describe the relationship between ISO~22000:2018, the ISO/TS~22002
        prerequisite programmes and the additional FSSC requirements.
  \item Interpret the key clauses and terminology of the scheme.
  \item Outline the roles of top management, the food-safety team and HACCP.
  \item Describe the certification, surveillance and re-certification cycle.
\end{itemize}

\coursesection{Course Modules}
\renewcommand{\arraystretch}{1.15}
\begin{tabular}{@{}p{0.9cm}p{9.4cm}p{1.9cm}@{}}
  \textbf{\#} & \textbf{Topic} & \textbf{Duration} \\[2pt]
  \hline
  1 & Introduction to food safety \& the FSSC~22000 scheme      & 60~min \\
  2 & ISO~22000:2018 — requirements and structure               & 90~min \\
  3 & Prerequisite programmes (ISO/TS~22002-1)                   & 75~min \\
  4 & HACCP principles \& hazard analysis                       & 90~min \\
  5 & Additional FSSC~22000 requirements                        & 60~min \\
  6 & Certification process, audits \& continual improvement     & 60~min \\
  7 & Review, examination \& feedback                            & 45~min \\
\end{tabular}

\coursesection{Assessment \& Certification}
Participants who attend the full programme and achieve the required pass mark
in the final examination are awarded this certificate of completion. The
certificate confirms attendance and successful assessment of the training
content described above; it is not, by itself, an accredited management-system
certification of the organisation.
